\section{Ventajas y limitaciones del método}

\subsection{Comparación exhaustiva de los métodos de guía en tiempo de inferencia}

\begin{center}
\begin{tabular}{p{3cm}p{3.5cm}p{3.5cm}p{3.5cm}}
\toprule
\textbf{Método} & \textbf{Manipulación de atención} & \textbf{DPS (problema inverso)} & \textbf{Guía por función de recompensa} \\
\midrule
Fundamento teórico & Ninguno (intuición de ingeniería) & Teorema de Bayes & Optimización con RL / restricciones KL \\
Condiciones de aplicación & Arquitecturas específicas & Mapeo $\mathcal{A}$ conocido & Función de recompensa diferenciable \\
Generalidad & Baja & Media & Alta \\
Dificultad de implementación & Media & Baja & Baja \\
Requisitos de entrenamiento & Ninguno & Ninguno & Ninguno (o mínimos) \\
Aplicaciones típicas & Control de composición & Superresolución / eliminación de ruido / reparación & Guía de preferencias arbitrarias \\
\bottomrule
\end{tabular}
\end{center}

\subsection{Ventajas}

\begin{enumerate}
    \item \textbf{Generalidad}: Siempre que se pueda definir una función de recompensa diferenciable, es aplicable. Funciona con múltiples modalidades como imágenes, video, movimiento y moléculas.
    \item \textbf{Sin entrenamiento adicional}: Aprovecha modelos preentrenados existentes, sin necesidad de recopilar nuevos datos ni realizar entrenamiento suplementario.
    \item \textbf{Implementación sencilla}: El código central solo requiere añadir un paso de guía por gradiente en el pipeline de generación.
    \item \textbf{Combinación flexible}: Se pueden aplicar múltiples funciones de recompensa simultáneamente para lograr una guía multicondicional.
    \item \textbf{Amigable computacionalmente}: En comparación con entrenar desde cero, la guía en tiempo de inferencia requiere menos recursos computacionales.
\end{enumerate}

\subsection{Limitaciones}

\begin{warningbox}{La función de recompensa debe ser diferenciable}
La guía en tiempo de inferencia depende de la propagación hacia atrás de gradientes; por lo tanto, la función de recompensa debe ser diferenciable. Para métricas de evaluación no diferenciables (como FID, IS, etc.), no pueden usarse directamente como señales de guía. Esta es una limitación fundamental del método actual.
\end{warningbox}

\begin{warningbox}{El equilibrio entre realismo y cumplimiento de condiciones es inevitable}
Aumentar la intensidad de la guía hace que las muestras intermedias se desvíen de la distribución de entrenamiento, lo que provoca que la red predictora de ruido genere predicciones inexactas en regiones fuera de distribución. El resultado final es:
\begin{itemize}
    \item Intensidad de guía baja: La salida es realista pero no cumple las condiciones.
    \item Intensidad de guía alta: Se cumplen las condiciones, pero la salida pierde realismo.
\end{itemize}
Buscar el peso de guía óptimo generalmente requiere ajuste empírico de hiperparámetros.
\end{warningbox}

Otras limitaciones incluyen:
\begin{itemize}
    \item Velocidad de inferencia más lenta: Más pasos temporales ofrecen un efecto de guía mejorado, pero incrementan el tiempo de inferencia.
    \item El aproximado de Tweedie tiene mayores errores en pasos de alto ruido: La calidad de la estimación de $\hat{x}_0$ en los primeros pasos temporales es pobre.
    \item No hay garantía teórica de convergencia al óptimo global.
\end{itemize}

\begin{knowledgebox}{Comparación con los métodos en tiempo de entrenamiento}
La guía en tiempo de inferencia complementa a los métodos en tiempo de entrenamiento (como ControlNet, modelos difusivos condicionales y ajuste fino RLHF):
\begin{itemize}
    \item \textbf{Métodos en tiempo de entrenamiento}: Ofrecen mayor estabilidad y velocidad de generación, pero requieren grandes cantidades de datos y recursos computacionales para el entrenamiento, además de que cada condición debe entrenarse por separado.
    \item \textbf{Métodos en tiempo de inferencia}: No requieren entrenamiento, tienen alta flexibilidad y permiten combinar múltiples condiciones al instante, aunque la estabilidad es menor y la velocidad de generación más lenta.
\end{itemize}
En la ingeniería práctica, la estrategia óptima suele ser una combinación de ambos: usar los métodos en tiempo de entrenamiento para tipos de condiciones comunes y fijos, y la guía en tiempo de inferencia para necesidades temporales y diversas.
\end{knowledgebox}

\subsection{Resumen de la sección}

Los métodos de guía en tiempo de inferencia presentan ventajas como una fuerte generalidad, la ausencia de necesidad de entrenamiento y una implementación sencilla, pero están limitados por el requisito de que la función de recompensa sea diferenciable y por el equilibrio inherente entre realismo y cumplimiento de condiciones.

